\section{Blind-Issuance Construction}
\label{sec:blind-signature}

This section presents the blind-issuance protocol that underlies ARC-SPRUCE.  The construction combines the
Ajtai commitment of Section~\ref{subsec:prelim-commit}, the BB-tran rational hash of
Section~\ref{subsec:prelim-vsis}, and the SmallWood NIZK-ARK of Section~\ref{subsec:prelim-nizk}.  The
signed message is the full vector $\vecmu=\vecm\Vert\veck$.

\subsection{Randomizing the BB-tran witness}
The holder first commits to
\[
  (\vecmu,r_0^H,r_1^H,\bar r),
\]
where $r_0^H$ and $r_1^H$ are holder-side contributions to the final BB-tran randomness and $\bar r$ is
commitment-only randomness.  The centered value $r_0\in\Rq^{\ell_{r_0}}$ is the long target-randomizer on the
linear side of the BB-tran target; its width is chosen so that Theorem~\ref{thm:htran-target-hiding} applies.
The issuer then samples bounded randomness
\[
  r_0^I,\ r_1^I
\]
and sends it to the holder.  The holder derives the final BB-tran randomness by coefficient-wise centering:
\begin{equation}
  r_0=\centerfunc(r_0^H+r_0^I),
  \qquad
  r_1=\centerfunc(r_1^H+r_1^I).
  \label{eq:bs-centered-randomness}
\end{equation}
Whenever $B_3-\mathbf 1_n r_1\in\Rq^{\times n}$, the holder computes the inverse witness
\[
  Z=(B_3-\mathbf 1_n r_1)^{-1}
\]
and the public BB-tran target
\begin{equation}
  T=B_0+B_1\vecmu+B_2r_0+Z
  =h_{\mathrm{tran},\vecmu,(r_0,r_1)}(B).
  \label{eq:bs-public-target}
\end{equation}
The pre-sign proof therefore binds the public target to the same hidden message and holder randomness already
fixed by the commitment, while the issuer controls only the additive randomness through
$(r_0^I,r_1^I)$.

\subsection{Protocol relations}
We record the two relations used later in the issuance and ARC layers.

\paragraph{Pre-sign issuance relation.}
The public statement is
\[
  x_{\mathrm{IssueBS}}=(\ComPublicParam,\vecc,r_0^I,r_1^I,T,B),
\]
and the hidden witness is
\[
  w_{\mathrm{IssueBS}}=(\vecmu,r_0^H,r_1^H,\bar r,r_0,r_1,Z).
\]
Define
\[
  \mathcal R_{\mathrm{IssueBS}}
  =
  \Bigl\{(x_{\mathrm{IssueBS}},w_{\mathrm{IssueBS}}):
  \text{all of the following hold}\Bigr\}
\]
where the constraints are
\begin{align}
  \vecc &= \mACom[\vecmu\Vert r_0^H\Vert r_1^H\Vert \bar r]^\top, \label{eq:issuebs-rel-commit}\\
  r_0 &= \centerfunc(r_0^H+r_0^I), \label{eq:issuebs-rel-center0}\\
  r_1 &= \centerfunc(r_1^H+r_1^I), \label{eq:issuebs-rel-center1}\\
  (B_3-\mathbf 1_n r_1)\Had Z &= 1, \label{eq:issuebs-rel-inverse}\\
  T &= B_0+B_1\vecmu+B_2r_0+Z, \label{eq:issuebs-rel-target}
\end{align}
and
\[
  \|\vecmu\|_\infty,\ \|r_0^H\|_\infty,\ \|r_1^H\|_\infty,\ \|\bar r\|_\infty,\ \|r_0\|_\infty,\ \|r_1\|_\infty
  \le \BoundMsg.
\]
Thus the proof certifies that the same hidden message and holder randomness both open the commitment and
give rise, after combining with the issuer randomness, to the public BB-tran target.

\paragraph{Post-issuance signature relation.}
For later knowledge statements, it is convenient to isolate the raw signature relation.  The public statement is
\[
  x_{\mathrm{SigBS}}=(\vecmu,\mA,B,\BoundMsg,\BoundSig),
\]
and the hidden witness is
\[
  w_{\mathrm{SigBS}}=(\vecu,r_0,r_1,Z).
\]
Define
\[
  \mathcal R_{\mathrm{SigBS}}
  =
  \Bigl\{(x_{\mathrm{SigBS}},w_{\mathrm{SigBS}}):
  (B_3-\mathbf 1_n r_1)\Had Z=1,
  \ \mA\vecu=B_0+B_1\vecmu+B_2r_0+Z,
  \ \|\vecu\|_\infty\le\BoundSig,
  \ \|\vecmu\|_\infty,\|r_0\|_\infty,\|r_1\|_\infty\le\BoundMsg
  \Bigr\}.
\]
Equivalently, this states that
\[
  \mA\vecu=h_{\mathrm{tran},\vecmu,(r_0,r_1)}(B)
\]
with all hidden values lying in the accepted bounded domains.

\subsection{Blind-issuance protocol and raw verification}
\label{subsec:bs-protocol}
We now record the shared-randomness blind-issuance protocol together with a raw post-issuance verification algorithm for the extracted witness relation.

\begin{description}
\item[$\Setup(\SecParam)$:]~
\begin{enumerate}[leftmargin=1.25em]
  \item Run
  \[
    \ComPublicParam\gets \ComScheme.\ComSetup(\SecParam,\PublicParamsGen)
  \]
  to obtain the Ajtai commitment key.
  \item Run
  \[
    \PublicParamsTrapdoor\gets \TDRH.\Setup(\SecParam,\PublicParamsGen)
  \]
  to obtain the BB-tran hash description $B$, the rational-hash randomness space $\mathcal X$, the admissibility
  error bound $\delta$, and the shortness bound $\BoundSig$.
  \item Output
  \[
    \PublicParamsBS=(\ComPublicParam,\PublicParamsTrapdoor,\BoundMsg,\BoundSig).
  \]
\end{enumerate}

\item[$\IKeyGen(\PublicParamsBS)$:]~
\begin{enumerate}[leftmargin=1.25em]
  \item Compute
  \[
    (\mA,\trapdoor)\gets \TDRH.\TrapGen(\PublicParamsTrapdoor).
  \]
  \item Output the issuer verification and signing keys
  \[
    \vk=\mA,
    \qquad
    \sk=\trapdoor.
  \]
\end{enumerate}

\item[$\Issue$:] The issuer $\Issuer$ holds $\sk$ while the holder $\Holder$ has input message $\vecmu$ and
issuer key $\vk$.
\begin{enumerate}[leftmargin=1.25em]
  \item The holder samples holder-side randomness
  \[
    r_0^H,\ r_1^H,\ \bar r
  \]
  coefficient-wise in $[-\BoundMsg,\BoundMsg]$, computes the commitment
  \[
    \vecc=\mACom[\vecmu\Vert r_0^H\Vert r_1^H\Vert \bar r]^\top,
  \]
  and sends $\vecc$ to the issuer.
  \item The issuer samples issuer-side randomness
  \[
    r_0^I,\ r_1^I
  \]
  coefficient-wise in $[-\BoundMsg,\BoundMsg]$ and sends $(r_0^I,r_1^I)$ to the holder.
  \item The holder computes the centered values $r_0,r_1$ as in
  \eqref{eq:bs-centered-randomness}.  If $B_3-\mathbf 1_n r_1\notin\Rq^{\times n}$, the holder aborts the current issuance and restarts from the issuer-randomness round with fresh issuer randomness; this failure probability is accounted for as part of the rational-hash admissibility error.  Otherwise it sets
  \[
    Z=(B_3-\mathbf 1_n r_1)^{-1},
    \qquad
    T=B_0+B_1\vecmu+B_2r_0+Z.
  \]
  \item The holder computes the pre-sign NIZK
  \[
    \nizkIssue\gets \nizkscheme.\ZKProve
      \bigl(w_{\mathrm{IssueBS}}\mid \mathcal R_{\mathrm{IssueBS}},x_{\mathrm{IssueBS}}\bigr)
  \]
  and sends $(T,\nizkIssue)$ to the issuer.
  \item The issuer verifies $\nizkIssue$ against $\mathcal R_{\mathrm{IssueBS}}$.  If verification fails, the issuer
  aborts.  Otherwise it samples
  \[
    \vecu\gets \TDRH.\SampPre(\trapdoor,T)
  \]
  and sends $\vecu$ to the holder.
  \item The holder checks
  \[
    \mA\vecu=T
    \qquad\text{and}\qquad
    \|\vecu\|_\infty\le \BoundSig.
  \]
  If either check fails, it aborts.  Otherwise it outputs the raw post-issuance witness
  \[
    \sig=(\vecu,r_0,r_1).
  \]
\end{enumerate}

\item[$\Verify(\vk,\vecmu,\sig)$:] On input
\[
  \sig=(\vecu,r_0,r_1),
\]
verification rejects if any of
\[
  \|\vecu\|_\infty>\BoundSig,
  \qquad
  \|\vecmu\|_\infty>\BoundMsg,
  \qquad
  \|r_0\|_\infty>\BoundMsg,
  \qquad
  \|r_1\|_\infty>\BoundMsg,
\]
fails.  It also rejects if $B_3-\mathbf 1_n r_1\notin\Rq^{\times n}$.  Otherwise it computes
\[
  Z=(B_3-\mathbf 1_n r_1)^{-1}
\]
and accepts iff
\[
  \mA\vecu=B_0+B_1\vecmu+B_2r_0+Z.
\]
\end{description}

The output $\sig=(\vecu,r_0,r_1)$ is a raw post-issuance witness for the BB-tran signature relation.  The
ARC layer uses it only inside later zero-knowledge showing proofs; the current section does not identify it
with an unlinkably rerandomized completed blind signature.

\subsection{Correctness and issuer-view hiding}
\begin{proposition}[Correctness]\label{prop:bs-correctness-prop}
Assume correctness of the NIZK and correctness of the trapdoor sampler.  Then honest execution of
$\Issue$ outputs a raw post-issuance witness accepted by $\Verify$ except with negligible probability.
\end{proposition}

\begin{proof}
By construction, the holder computes $r_0,r_1,Z,T$ satisfying
\[
  (B_3-\mathbf 1_n r_1)\Had Z=1,
  \qquad
  T=B_0+B_1\vecmu+B_2r_0+Z.
\]
Completeness of the pre-sign NIZK makes the issuer accept, and correctness of $\SampPre$ yields
$\mA\vecu=T$ with $\|\vecu\|_\infty\le \BoundSig$ except with negligible failure probability.  The final
verification algorithm recomputes the same $Z$ and therefore accepts.
\end{proof}

\begin{proposition}[Honest-issuer public-surface hiding]\label{prop:bs-issuer-view-hiding}
Assume:
\begin{enumerate}[leftmargin=1.25em]
  \item the issuer samples $r_0^I,r_1^I$ honestly and independently from the prescribed bounded uniform
  distributions;
  \item the pre-sign NIZK is zero knowledge;
  \item the Ajtai commitment is hiding as in Lemma~\ref{lem:ajtai-hiding};
  \item the centering map satisfies Lemma~\ref{lem:center-uniform}; and
  \item Theorem~\ref{thm:htran-target-hiding} holds for the protocol distribution of $r_0$.
\end{enumerate}
Then the public issuance surface
\[
  (\vecc,r_0^I,r_1^I,T,\nizkIssue)
\]
is computationally hiding with respect to the holder message $\vecmu$, up to the NIZK simulation error, the
commitment-hiding error, the target-hiding error, and the stated admissibility/restart error.
\end{proposition}

\begin{proof}
Let $\vecmu_0,\vecmu_1$ be two challenge messages and let $b\gets\{0,1\}$ be the hidden bit.
We consider an honest issuance to $\vecmu_b$.

First replace the real pre-sign proof $\nizkIssue$ by a simulated proof for the same public statement.  By zero
knowledge, this changes the issuer's distinguishing advantage by at most the NIZK simulation error.  After this
step, the proof transcript is no longer a possible leakage channel, so the relevant public view is reduced to
\[
  (\vecc,r_0^I,r_1^I,T).
\]

Second, the Ajtai commitment $\vecc$ hides the committed vector
\[
  (\vecmu_b,r_0^H,r_1^H,\bar r)
\]
by Lemma~\ref{lem:ajtai-hiding}.  In particular, it hides both the challenge message and the holder-side
randomness that will later be combined with the issuer randomness.

Third, the issuer contributions $r_0^I,r_1^I$ are sampled independently of the challenge bit.  For every fixed
holder-side contribution, Lemma~\ref{lem:center-uniform} implies that
\[
  r_0=\centerfunc(r_0^H+r_0^I)
\]
has exactly the coefficient-wise bounded uniform distribution required by Lemma~\ref{lem:htran-lhl}.  If the
protocol restarts until $B_3-\mathbf 1_n r_1$ is admissible, the conditioning depends only on $r_1$ and does
not change the bounded-uniform distribution of the independently sampled $r_0$.

For honestly generated public parameters, and in particular for uniformly sampled $B_2$,
Theorem~\ref{thm:htran-target-hiding} implies that for every fixed admissible value of $r_1$ the public target
\[
  T=B_0+B_1\vecmu_b+B_2r_0+(B_3-\mathbf 1_n r_1)^{-1}
\]
statistically hides $\vecmu_b$.  Averaging this pointwise bound over the accepted distribution of $r_1$ preserves
the same bound.  Since $r_0^I$ and $r_1^I$ are independent of $b$, combining commitment hiding, independent
issuer randomness, and the LHL-based target-hiding theorem proves the stated public-surface hiding claim.
\end{proof}

\begin{remark}[Security status]\label{rem:bs-security-status}
The non-blind signature layer has the conditional DKLW guarantees recorded in
Section~\ref{subsec:prelim-signatures}.  The blind-issuance protocol established in this section provides
correctness and honest-issuer public-surface hiding.  A separate reduction or assumption is still needed for
one-more unforgeability of the
interactive blind-issuance layer itself.
\end{remark}
